% !TEX root = task-1-17-3.tex
% Editable task bank. Append new BankTask blocks at the end of the document.
\documentclass[12pt,a4paper]{article}
\usepackage{physics-textbook}
\usepackage{graphicx,needspace}
\usetikzlibrary{decorations.pathmorphing}
\newif\ifShowSolutions
\ifdefined\TutorEdition
  \ShowSolutionstrue
\else
  \ShowSolutionsfalse % Change to \ShowSolutionstrue to display solutions.
\fi
\newcounter{banktask}
\newsavebox{\BankTaskBox}
\newcommand{\BankTaskContent}[3]{%
  \noindent\textbf{\thebanktask.}\quad #2\par
  \ifShowSolutions
    \Needspace{50mm}% Keep the solution heading with its opening text or diagram.
    #3\par
    \ifdefined\TutorEdition
      \ifcsname TutorSource#1\endcsname
        \nopagebreak\smallskip{\small\csname TutorSource#1\endcsname\par}%
      \fi
    \fi
  \fi
}
\newcommand{\BankTask}[3]{%
  \refstepcounter{banktask}\par\medskip
  \sbox{\BankTaskBox}{\begin{minipage}{\linewidth}%
    \BankTaskContent{#1}{#2}{#3}\end{minipage}}%
  \ifdim\dimexpr\ht\BankTaskBox+\dp\BankTaskBox\relax>\dimexpr\textheight-20mm\relax
    \Needspace{35mm}\BankTaskContent{#1}{#2}{#3}%
  \else
    \Needspace{\dimexpr\ht\BankTaskBox+\dp\BankTaskBox\relax}%
    \noindent\usebox{\BankTaskBox}\par
  \fi
  \medskip
}
\begin{document}
\BookSubsection{1.17.3}{Пружинный маятник}


% Source sdamgia-717: https://phys-ege.sdamgia.ru/problem?id=717
\BankTask{sdamgia-717}{
Гиря массой $4\,\text{кг}$, подвешенная на стальной пружине, совершает свободные колебания с периодом $2\,\text{с}$. С каким периодом будет совершать свободные колебания гиря массой $1\,\text{кг}$, подвешенная на этой пружине?
}{
\smallskip\noindent\textbf{Решение}\par
Период колебаний пружинного маятника связан с массой груза и жесткостью пружины соотношением:

$T=2 \pi \sqrt{\frac{m}{k}} .$

Следовательно, если оставить пружину прежней, а массу груза уменьшить в $4$ раза, период уменьшится в $2$ раза и будет равен $1\,\text{с}$.

Ответ: $1$.
}

% Source sdamgia-636: https://phys-ege.sdamgia.ru/problem?id=636
\BankTask{sdamgia-636}{
Груз, подвешенный на пружине жесткостью $400\,\text{Н/м},$ совершает свободные гармонические колебания. Какой должна быть жёсткость пружины, чтобы частота колебаний этого груза увеличилась в $2$ раза?
}{
\smallskip\noindent\textbf{Решение}\par
Частота колебаний пружинного маятника связана с жесткостью пружины и массой груза соотношением:

$\nu = \frac{1}{2 \pi} \sqrt{\frac{k}{m}} .$

Следовательно, при неизменной массе груза для увеличения частоты колебаний в два раза, необходимо увеличить жесткость пружины в $4$ раза. Таким образом, жесткость пружины должна быть равна:

$4 \cdot 400\,\text{Н/м}=1600\,\text{Н/м}.$

Ответ: $1600$.
}

% Source sdamgia-11786: https://phys-ege.sdamgia.ru/problem?id=11786
\BankTask{sdamgia-11786}{
Шарик массой $0{,}4\,\text{кг}$, подвешенный на пружине, совершает свободные гармонические колебания вдоль вертикальной прямой. Какой должна быть масса шарика, чтобы период его свободных вертикальных гармонических колебаний на этой же пружине был в $2$ раза меньше?
}{
\smallskip\noindent\textbf{Решение}\par
Период колебаний пружинного маятника выражается формулой: $T = 2 \pi \sqrt{\frac{m}{k}} .$ Следовательно, для уменьшения периода колебаний в два раза необходимо уменьшить массу шарика в $4$ раза, то есть масса шарика должна быть равной $0{,}1\,\text{кг}$.

Ответ: $0{,}1$.
}

% Source sdamgia-10462: https://phys-ege.sdamgia.ru/problem?id=10462
\BankTask{sdamgia-10462}{
Период свободных колебаний пружинного маятника равен $0{,}5\,\text{с}$. Каким станет период свободных колебаний этого маятника, если массу груза маятника увеличить в $2$ раза, а жесткость пружины вдвое уменьшить?
}{
\smallskip\noindent\textbf{Решение}\par
Период свободных колебаний пружинного маятника связан с жесткостью пружины и массой груза соотношением $T=2 \pi \sqrt{\frac{m}{k}} .$ При увеличении массы груза в $2$ раза и уменьшении жесткости пружины вдвое период колебаний увеличится в $2$ раза, то есть станет равным $1\,\text{с}$.

Ответ: $1$.
}

% Source sdamgia-7281: https://phys-ege.sdamgia.ru/problem?id=7281
\BankTask{sdamgia-7281}{
Груз на длинной легкой пружине совершает колебания с частотой $0{,}5\,\text{Гц}$. Пружину разрезали на $4$ равные части и прикрепили к одной из частей тот же груз. Чему стал равен период колебаний получившегося пружинного маятника?
}{
\smallskip\noindent\textbf{Решение}\par
Период колебаний потенциальной энергии пружинного маятника пропорционален периоду колебаний груза, который определяется выражением:

$T=2 \pi \sqrt{\frac{m}{k}} .$

При последовательном соединении пружин общая жесткость получается из формулы: $1/k= ( 1 / k_1 + 1 / k_2 + 1 / k_3 + ... + 1 / k_n ) .$ Отсюда следует, что жесткость $1/4$ части исходной пружины равна $4k$. Увеличение жесткости пружины в $4$ раза приводит к уменьшению периода колебаний в два раза, а частота, являясь обратной периоду величиной, в свою очередь должна увеличиться в два раза и станет равной $1\,\text{Гц}$. Тогда период колебаний получившегося маятника будет равен $1\,\text{с}$.

Ответ: $1$.
}

% Source sdamgia-29738: https://phys-ege.sdamgia.ru/problem?id=29738
\BankTask{sdamgia-29738}{
Груз массой $2\,\text{кг}$, закрепленный на пружине жесткостью $200\,\text{Н/м}$, совершает гармонические колебания с амплитудой $10\,\text{см}$. Какова максимальная скорость груза?
}{
\smallskip\noindent\textbf{Решение}\par
При колебаниях пружинного маятника выполняется закон сохранения энергии. Так как в положении равновесия полная энергия равна максимальной кинетической энергии груза 

$E_k= \frac{m v _{\max} ^2}{2} ,$

 а в крайних точках полная энергия равна максимальной потенциальной энергии пружины $E_p= \frac{k A ^2}{2} ,$ то: 

$\frac{m v _{\max} ^2}{2} = \frac{k A ^2}{2} .$

Откуда максимальная скорость груза: 

$v _{\max} =A \sqrt{\frac{k}{m}} =0{,}1 \cdot \sqrt{\frac{200}{2}} =1 \,\text{м/с}.$ 

Ответ: $1\,\text{м/с}$.
}

% Source sdamgia-3720: https://phys-ege.sdamgia.ru/problem?id=3720
\BankTask{sdamgia-3720}{
\begin{center}\begin{tikzpicture}[task graph,x=.75cm,y=.32cm]
\draw[task grid,xstep=.5,ystep=1] (0,0) grid (4.5,9);\draw[task axis] (0,0)--(4.9,0) node[right] {$x$, см};\draw[task axis] (0,0)--(0,9.6) node[above] {$F$, Н};\draw[task curve] (0,0)--(4.5,9);
\foreach \x in {1,2,3,4}{\node[below] at (\x,0) {$\x$};}\foreach \y in {2,4,6}{\node[left] at (0,\y) {$\y$};}\node[below left] at (0,0) {$0$};
\end{tikzpicture}\end{center}На графике приведена зависимость модуля силы упругости $F$ растянутой пружины от величины ее растяжения $x.$ Найти период свободных колебаний груза массой $0{,}5\,\text{кг}$, подвешенного на этой пружине.
}{
\smallskip\noindent\textbf{Решение}\par
Определим жесткость пружины. Согласно закону Гука, сила упругости пропорциональна удлинению пружины $F=kx,$ что и подтверждает данный график. Используя любую точку на графике для жесткости пружины, имеем: 

$k= \frac{F}{x} = \frac{4\,\text{Н}}{0{,}02\,\text{м}} =200\,\text{Н/м}.$

Период свободных колебаний пружинного маятника связан с жесткостью пружины и массой груза соотношением:

$T=2 \pi \sqrt{\frac{m}{k}} =2 \cdot 3{,}14 \cdot \sqrt{\frac{0{,}5\,\text{кг}}{200\,\text{Н/м}}} =0{,}314\,\text{с}\approx0{,}3\,\text{с}.$

Ответ: $0{,}3\,\text{с}$.
}

% Source sdamgia-50848: https://phys-ege.sdamgia.ru/problem?id=50848
\BankTask{sdamgia-50848}{
\begin{center}\begin{tikzpicture}[task graph]
\draw[task curve] (0,0)--(5,0);\draw[task curve] (0,0)--(0,2);
\foreach \x in {.1,.35,...,4.9}{\draw[task tick] (\x,0)--++(-.18,-.18);}\foreach \y in {.15,.4,...,1.9}{\draw[task tick] (0,\y)--++(-.18,-.18);}
\draw[task curve] (0,1)--(.5,1)--(.65,1.65)--(.95,.4)--(1.25,1.65)--(1.55,.4)--(1.85,1.65)--(2.15,.4)--(2.45,1.65)--(2.75,.4)--(2.9,1)--(3.1,1);
\draw[task curve,fill=orange!25] (3.1,.55) rectangle (4.8,1.7);\draw[task curve,fill=white] (3.45,.25) circle (.25);\draw[task curve,fill=white] (4.45,.25) circle (.25);
\end{tikzpicture}\end{center}Тележка массой $0{,}5\,\text{кг}$, прикрепленная к горизонтальной пружине жёсткостью $200\,\text{Н/м}$, совершает свободные гармонические колебания (см. рис.). Максимальная скорость тележки равна $3\,\text{м}/\text{с}$. Какова амплитуда колебаний тележки? Массой колёс можно пренебречь.
}{
\smallskip\noindent\textbf{Решение}\par
При гармонических колебаниях выполняется закон сохранения энергии: $\frac{kA ^2}{2} = \frac{m v _{\max} ^2}{2} ,$ откуда найдем амплитуду колебаний тележки:

$A = v _{\max} \sqrt{\frac{m}{k}} = 3 \cdot \sqrt{\frac{0{,}5}{200}} = 0{,}15\,\text{м} = 15\,\text{см}.$

Ответ: $15\,\text{см}$.
}


% Source B86809: https://ege.fipi.ru/bank/questions.php?search=1&proj=BA1F39653304A5B041B656915DC36B38&qid=B86809
\BankTask{B86809}{
Груз, подвешенный на лёгкой пружине жёсткостью $400\,\text{Н}/\text{м}$, совершает свободные вертикальные гармонические колебания. Пружину какой жёсткости надо взять вместо первой пружины, чтобы период свободных колебаний этого груза стал в $2$ раза меньше? Результат выразите в $\text{Н}/\text{м}$.
}{
\smallskip\noindent\textbf{Решение от GPT}\par
$T=2\pi\sqrt{m/k}$. При $T_2=T_1/2$ и той же массе $k_2=4k_1=1600\,\text{Н}/\text{м}$.
}

% Source 6DEE76: https://ege.fipi.ru/bank/questions.php?search=1&proj=BA1F39653304A5B041B656915DC36B38&qid=6DEE76
\BankTask{6DEE76}{
Период свободных колебаний пружинного маятника равен $0{,}5\,\text{с}$. Каким станет период свободных колебаний этого маятника, если массу груза маятника увеличить в $2$ раза, а жёсткость пружины вдвое уменьшить? Результат выразите в $\text{с}$.
}{
\smallskip\noindent\textbf{Решение от GPT}\par
$T_2/T_1=\sqrt{(m_2/m_1)(k_1/k_2)}=\sqrt{2\cdot2}=2$. Значит, $T_2=1\,\text{с}$.
}

% Source 1B36B0: https://ege.fipi.ru/bank/questions.php?search=1&proj=BA1F39653304A5B041B656915DC36B38&qid=1B36B0
\BankTask{1B36B0}{
Шарик массой $0{,}4\,\text{кг}$, подвешенный на пружине, совершает свободные гармонические колебания вдоль вертикальной прямой. Какой должна быть масса шарика, чтобы период его свободных вертикальных гармонических колебаний на этой же пружине был в $2$ раза меньше? Результат выразите в $\text{кг}$.
}{
\smallskip\noindent\textbf{Решение от GPT}\par
Для одной пружины $T^2\sim m$. Поэтому $m_2=m_1(T_2/T_1)^2=0{,}4/4=0{,}1\,\text{кг}$.
}

% Source 1BF6BB: https://ege.fipi.ru/bank/questions.php?search=1&proj=BA1F39653304A5B041B656915DC36B38&qid=1BF6BB
\BankTask{1BF6BB}{
Частота свободных гармонических колебаний пружинного маятника равна $4\,\text{Гц}$. Какой будет частота свободных колебаний этого маятника, если массу его груза увеличить в $4$ раза? Результат выразите в $\text{Гц}$.
}{
\smallskip\noindent\textbf{Решение от GPT}\par
$\nu\sim m^{-1/2}$ при неизменной пружине. Тогда $\nu_2=\nu_1/\sqrt4=2\,\text{Гц}$.
}

% Source eeB001: https://ege.fipi.ru/bank/questions.php?search=1&proj=BA1F39653304A5B041B656915DC36B38&qid=eeB001
\BankTask{eeB001}{
В таблице представлены данные о положении грузика, прикреплённого к пружине и совершающего гармонические колебания вдоль горизонтальной оси $Ox$, в различные моменты времени.
\begin{center}{\small\setlength{\tabcolsep}{3pt}\begin{tabular}{|c|c|c|c|c|c|c|c|c|c|c|c|c|c|c|c|c|}\hline $t,\,\text{с}$ & $0{,}0$ & $0{,}1$ & $0{,}2$ & $0{,}3$ & $0{,}4$ & $0{,}5$ & $0{,}6$ & $0{,}7$ & $0{,}8$ & $0{,}9$ & $1{,}0$ & $1{,}1$ & $1{,}2$ & $1{,}3$ & $1{,}4$ & $1{,}5$ \\\hline $x,\,\text{мм}$ & $0$ & $2{,}5$ & $4{,}5$ & $6{,}0$ & $7{,}0$ & $7{,}5$ & $7{,}0$ & $6{,}0$ & $4{,}5$ & $2{,}5$ & $0$ & $-2{,}5$ & $-4{,}5$ & $-6{,}0$ & $-7{,}0$ & $-7{,}5$ \\\hline\end{tabular}}\end{center}
\begin{enumerate}
\item Период колебаний грузика равен $1\,\text{с}$.
\item Амплитуда колебаний грузика равна $15\,\text{мм}$.
\item Потенциальная энергия пружины маятника в момент времени $1{,}5\,\text{с}$ минимальна.
\item Кинетическая энергия грузика в момент времени $1{,}0\,\text{с}$ максимальна.
\item Модуль ускорения грузика в момент времени $1{,}5\,\text{с}$ максимален.
\end{enumerate}
}{
\smallskip\noindent\textbf{Решение от GPT}\par
Из таблицы $A=7{,}5\,\text{мм}$, $T=2\,\text{с}$. При $t=1{,}0\,\text{с}$ $x=0$, поэтому кинетическая энергия максимальна; при $t=1{,}5\,\text{с}$ $|x|=A$, поэтому $|a|=\omega^2A$ максимально. Ответ: $45$.
}

% Source C1C6AE: https://ege.fipi.ru/bank/questions.php?search=1&proj=BA1F39653304A5B041B656915DC36B38&qid=C1C6AE
\BankTask{C1C6AE}{
На рисунке приведены зависимости координат $x$ двух грузов $1$ и $2$ от времени $t$ в процессе их малых колебаний на гладкой горизонтальной плоскости. Грузы имеют одинаковую массу $m$ и прикреплены к пружинам жёсткостью $k_{1}$ (груз $1$) и жёсткостью $k_{2}$ (груз $2$). Колебания груза $1$ обозначены сплошной линией, груза $2$ – пунктиром.
\begin{center}\begin{tikzpicture}[x=.96cm,y=.8cm,>=stealth]
\foreach \q in {0,1,2,3,4,5,6} \draw[thin,black!65] (\q,-2)--(\q,2);
\draw[thin,black!65] (0,2)--(6,2) (0,-2)--(6,-2);
\foreach \h in {-1.5,-1,-.5,.5,1,1.5} \draw[thin,black!65] (0,\h)--(1.1,\h) (1.4,\h)--(2,\h) (2.8,\h)--(3.5,\h) (4,\h)--(5,\h) (5.3,\h)--(6,\h);
\draw[->] (0,0)--(6.6,0) node[right] {$t,\,\text{с}$};
\draw[->] (0,-2)--(0,2.7) node[above] {$x,\,\text{см}$};
\foreach \x in {1,2,3,4,5,6} \node[below] at (\x,0) {$\x$};
\foreach \y in {-2,-1,0,1,2} \node[left] at (0,\y) {$\y$};
\draw[very thick,domain=0:6,samples=100] plot (\x,{2*sin(90*\x)});
\draw[thick,dashed,domain=0:6,samples=100] plot (\x,{cos(90*(\x-.25))});
\node at (5.8,1.85) {$1$};\node at (3.2,1.0) {$2$};
\end{tikzpicture}\end{center}
\begin{enumerate}
\item Периоды колебаний грузов одинаковы.
\item Максимальные потенциальные энергии деформированных пружин одинаковы.
\item Жёсткость пружины $k_{1}$ меньше жёсткости пружины $k_{2}$.
\item Амплитуды колебаний грузов одинаковы.
\item Максимальная кинетическая энергия груза $1$ в четыре раза больше, чем максимальная кинетическая энергия груза $2$.
\end{enumerate}
}{
\smallskip\noindent\textbf{Решение от GPT}\par
По графику периоды одинаковы, значит при одинаковых массах $k_1=k_2$. Амплитуды $A_1=2\,\text{см}$ и $A_2=1\,\text{см}$ различаются вдвое. Поэтому $E_{\text{к max}}=kA^2/2$ у первого груза вчетверо больше. Ответ: $15$.
}

% Source EC9E02: https://ege.fipi.ru/bank/questions.php?search=1&proj=BA1F39653304A5B041B656915DC36B38&qid=EC9E02
\BankTask{EC9E02}{
На гладком горизонтальном столе брусок массой $M$, прикреплённый к вертикальной стене лёгкой горизонтальной пружиной жёсткостью $k$, совершает свободные гармонические колебания с амплитудой $A$.
\TaskMatchingColumns[.53\linewidth]{Физические величины}{
А) частота колебаний бруска\par
Б) амплитуда ускорения бруска
}{Формулы}{
1) $\frac{1}{2\pi}\sqrt{\frac{M}{k}}$\par
2) $A\frac{k}{M}$\par
3) $\frac{1}{2\pi}\sqrt{\frac{k}{M}}$\par
4) $A\sqrt{\frac{k}{M}}$\par
}\begin{center}\begin{tabular}{|c|c|}\hline А & Б \\\hline\rule{0pt}{1.2em}\hspace{1cm} & \hspace{1cm} \\\hline\end{tabular}\end{center}
}{
\smallskip\noindent\textbf{Решение от GPT}\par
$\omega=\sqrt{k/M}$, поэтому $\nu=\omega/(2\pi)$ (формула $3$), $a_{\max}=\omega^2A=Ak/M$ (формула $2$). Ответ: $32$.
}

% Source DF724F: https://ege.fipi.ru/bank/questions.php?search=1&proj=BA1F39653304A5B041B656915DC36B38&qid=DF724F
\BankTask{DF724F}{
Подвешенный на пружине груз совершает свободные вертикальные гармонические колебания. Пружину заменили на другую, жёсткость которой больше, оставив массу груза и амплитуду колебаний неизменными. Как при этом изменятся частота свободных колебаний груза и его максимальная скорость?
\begin{enumerate}
\item увеличится
\item уменьшится
\item не изменится
\end{enumerate}
\TaskChangeTable{Частота свободных колебаний груза}{Максимальная скорость груза}
}{
\smallskip\noindent\textbf{Решение от GPT}\par
$\nu=(2\pi)^{-1}\sqrt{k/m}$, $v_{\max}=2\pi\nu A$. При увеличении $k$ и неизменных $m,A$ обе величины увеличиваются. Ответ: $11$.
}

% Source FA9CF6: https://ege.fipi.ru/bank/questions.php?search=1&proj=BA1F39653304A5B041B656915DC36B38&qid=FA9CF6
\BankTask{FA9CF6}{
Ученику необходимо на опыте обнаружить зависимость периода свободных колебаний пружинного маятника от массы груза. У него имеется пять пружинных маятников, характеристики которых приведены в таблице. Какие два маятника необходимо взять ученику, чтобы провести данное исследование? Грузы сплошные.
\begin{center}
\begin{tabular}{|c|c|c|c|}
\hline
\shortstack{№\\маятника} & \shortstack{Жёсткость пружины,\\\(\text{Н}/\text{м}\)} & \shortstack{Объём груза,\\$\text{см}^{3}$} & \shortstack{Материал, из которого\\сделан груз} \\
\hline
\(1\) & \(40\) & \(30\) & алюминий \\
\hline
\(2\) & \(60\) & \(60\) & алюминий \\
\hline
\(3\) & \(40\) & \(30\) & медь \\
\hline
\(4\) & \(10\) & \(30\) & алюминий \\
\hline
\(5\) & \(10\) & \(60\) & медь \\
\hline
\end{tabular}
\end{center}
}{
\smallskip\noindent\textbf{Решение от GPT}\par
Период пружинного маятника равен $T=2\pi\sqrt{m/k}$, поэтому нужно сохранить жёсткость пружины и изменить массу груза. У маятников 1 и 3 одинаковая жёсткость $k=40\,\text{Н/м}$, но разные массы $m=\rho V$, поскольку при одинаковом объёме материалы различаются. Ответ: 1 и 3. По приведённым данным также подходит пара 4 и 5: жёсткости одинаковы, а массы различны, поэтому выбор пары в условии не единственен.
}

% Source 50844: https://phys-ege.sdamgia.ru/problem?id=50844
\BankTask{50844}{
Груз массой $2\,\text{кг}$, закреплённый на пружине жесткостью $200\,\text{Н/м}$, совершает горизонтальные гармонические колебания. Максимальное ускорение груза при этом равно $10\,\text{м/с}^{2}$. Какова максимальная скорость груза?
}{
\smallskip\noindent\textbf{Решение}\par
При гармонических колебаниях максимальная сила $F_{\max}=kA,$ где A  — амплитуда колебаний. По второму закону Ньютона $F_{\max}=ma_{\max}.$ Тогда амплитуда колебаний равна

\[A = \frac{ma_{\max}}{k} .\]

Для гармонических колебаний выполняется закон сохранения энергии:

\[\frac{kA ^2}{2} = \frac{m v _{\max} ^2}{2} .\]

Объединяя уравнения, находим максимальную скорость груза:

\[v _{\max}=a_{\max} \sqrt{\frac{m}{k}} =10 \cdot \sqrt{\frac{2}{200}} =1 \,\text{м/с}.\]

Ответ: $1\,\text{м/с}$.
}

\end{document}
